\section{Ordinary primes}\label{sec:ordinary}

This section proves the lower bound $\Iprof_{\mathcal C}h^2$ for the contribution of the
ordinary primes $p\nmid2q(q^2-4\sigma)$ to $\sum_p\vG p(T)\log p$ (\cref{thm:ordinary}); it is the
largest term of $\aconst_\sigma(q)$ and is used in \cref{thm:arithmetic}. Primes $p<\delta h$ are
handled by the crude bound of \cref{prop:uniform} and primes $p>2K$ cost nothing
(\cref{lem:large-primes}). The work is in the range $p\asymp h$, where we bound $\vG p(T)$ by a
small optimisation problem over the residue classes modulo $p$ (\cref{subsec:two-bounds}), pass to
a limiting profile in $p/h$ (\cref{subsec:profile}), and integrate it against $\log p$.

Throughout, $p$ is ordinary, so $w$ and $1-w$ are $p$-adic units.

\begin{lemma}\label{lem:large-primes}
If $p>2K$ is ordinary, then $\vG p(T)\ge0$.
\end{lemma}

\begin{proof}
All factorials in $\nrm$ involve integers smaller than $p$, so $v_p(\nrm)=0$. The polynomial
$D_K$ is monic with $p$-integral coefficients, so the quotient of $u^{i+j+1}N_A(u)$ by $D_K$ has
$p$-integral coefficients and so has its derivative; by \cref{lem:moment-sums}(c) its moment sum
is $p$-integral. The residues $x_k^{i+j+1}N_A(x_k)/D_K'(x_k)$ are $p$-integral because
$D_K'(x_k)=\pm k!(K-1-k)!$, and $\tau_k$ involves only the odd numbers $2j+1<2K<p$.
\end{proof}

The interesting range is $p\asymp h$. Fix $\delta>0$. For $h$ large, every prime
$p\ge\delta h$ satisfies $p^2>3h>2K$ (both constructions have $2K<3h$), and we assume this from
now on. Then every difference
$l-k$ of pole indices and every odd number below $2K$ has $p$-adic valuation at most one.

\subsection{Residue classes and local functionals}\label{subsec:distribution}

\emph{Heuristic.} Modulo $p$ the poles fall into residue classes. In the coordinate $u=a+pz$
around a class $a$, the poles of that class are at distance $O(1)$ while all others are units,
so the weight $\varphi$ has a definite $p$-adic size on each class. Splitting the series over
the residue classes of $m$ (\cref{lem:distribution}) turns $\mu_X$ into a sum of local
functionals, one per class, which are integral (\cref{lem:local-integrality}).


For $a\in\{1,\dots,p\}$ let $\Pi_a=\{k<K:x_k\equiv a\pmod p\}$, $c_a=\abs{\Pi_a}$, and for
$k\in\Pi_a$ let $z_k=(x_k-a)/p\in\Z_p$; these are pairwise incongruent modulo $p$, because
$\abs{l-k}<K<p^2$. Let $\rest{\Z_p}$ be the ring of restricted power series
$\sum_jg_jz^j$ with $g_j\in\Z_p$, $g_j\to0$, and $\rest{\Q_p}=\rest{\Z_p}[1/p]$. Put $v=w^p$;
then $v\equiv w\pmod p$, so $v$ and $1-v$ are units. Let $M_r(v)=\sum_{n\ge0}n^rv^n\in\Z_p$ be
the rational moment ($M_0=1/(1-v)$, and $M_r=\sum_{i=1}^r S(r,i)\,i!\,v^i/(1-v)^{i+1}$ with
Stirling numbers $S(r,i)$), and define the continuous linear map
\[
  \mathsf M_v\colon\rest{\Q_p}\to\Q_p,\qquad \mathsf M_v\Bigl(\sum_jg_jz^j\Bigr)=\sum_jg_jM_j(v),
\]
which maps $\rest{\Z_p}$ into $\Z_p$. Let $V_a=\rest{\Q_p}\oplus\bigoplus_{k\in\Pi_a}\Q_p\cdot
(z-z_k)^{-1}$ and define $\ell_a\colon V_a\to\Q_p[X]$ by $\ell_a(G)=\mathsf M_v(G')$ for
$G\in\rest{\Q_p}$ and, for $k\in\Pi_a$,
\begin{equation}\label{eq:ell-pole}
  \ell_a\Bigl(\frac1{z-z_k}\Bigr)=\frac{p^2}{w^a}\Bigl[\mu_X\Bigl(\frac1{u-x_k}\Bigr)
  -\sum_{\substack{b=1\\b\ne a}}^{p}\frac{w^b}{p}\,
  \mathsf M_v\Bigl(\frac{\mathrm d}{\mathrm dz}\,\frac{1}{b+pz-x_k}\Bigr)\Bigr].
\end{equation}
For $R\in\Rsp$ let $R_a(z)=R(a+pz)$. If $R$ is a polynomial, $R_a\in\Q_p[z]$; if
$R=1/(u-x_k)$ then $R_a=p^{-1}(z-z_k)^{-1}$ for $k\in\Pi_a$, and otherwise $a-x_k$ is a unit and
$R_a\in\rest{\Z_p}$. Thus $R_a\in V_a$.

\begin{lemma}[distribution identity]\label{lem:distribution}
For every $R\in\Rsp$, $\displaystyle\mu_X(R)=\sum_{a=1}^{p}\frac{w^a}{p}\,\ell_a(R_a)$.
\end{lemma}

\begin{proof}
For a polynomial $g$ one has $\sum_{m\ge1}g(m)w^m=\sum_{a=1}^pw^a\sum_{n\ge0}g(a+pn)w^{pn}$ in
$\Q[[w]]$; both sides are rational functions of $w$, so the identity holds at our $w$, in
$\Q_p$. Apply it to $g=R'$ for a polynomial $R$, using $(R_a)'(z)=pR'(a+pz)$. For
$R=1/(u-x_k)$ with $k\in\Pi_{a_0}$ the identity is equivalent to the definition~\eqref{eq:ell-pole}
of $\ell_{a_0}(1/(z-z_k))$.
\end{proof}

\begin{lemma}[integrality]\label{lem:local-integrality}
Let $a\in\{1,\dots,p\}$ and $P_a(z)=\prod_{k\in\Pi_a}(z-z_k)$.
\begin{enumerate}[(a)]
\item For $k\in\Pi_a$, $\ell_a(1/(z-z_k))\in\Z_p[X]$.
\item If $N\in\rest{\Z_p}$, then $G=N/P_a\in V_a$ and $\ell_a(G)\in\Z_p[X]$.
\end{enumerate}
\end{lemma}

\begin{proof}
(a) In $p^2\mu_X(1/(u-x_k))=-4\sigma^ky^{2k+1}(p^2X-p^2\tau_k)$ the number $p^2\tau_k$ is
$p$-integral, since $v_p(2j+1)\le1$ for $2j+1<2K$. In the sum over $b\ne a$,
$p^{-1}\frac{\mathrm d}{\mathrm dz}(b+pz-x_k)^{-1}=-(b+pz-x_k)^{-2}$ lies in $\rest{\Z_p}$.
(b) Division of $z^j$ by the monic polynomial $P_a\in\Z_p[z]$ has integral quotient and
remainder, so $N=QP_a+N_0$ with $Q\in\rest{\Z_p}$ and $N_0\in\Z_p[z]$ of degree less than $c_a$.
Then $G=Q+\sum_{k\in\Pi_a}\frac{N_0(z_k)}{P_a'(z_k)}\frac1{z-z_k}$, and $P_a'(z_k)$ is a unit
because the $z_k$ are pairwise incongruent. Apply~(a).
\end{proof}

\emph{Local weights.} Let $Z_a=\{1\le j\le A:j\equiv a\pmod p\}$ and $n_a=\abs{Z_a}$. For $k\notin\Pi_a$ the factor
$a+pz-x_k$ is a unit of $\rest{\Z_p}$, and for $j\notin Z_a$ so is $a+pz-j$. Hence
\begin{equation}\label{eq:local-weight}
  \varphi(a+pz)=p^{\,2n_a-c_a}\,\frac{N_a(z)}{P_a(z)},\qquad
  N_a(z)=(a+pz)\prod_{j\in Z_a}\Bigl(z-\frac{j-a}p\Bigr)^2\,E_a(z),
\end{equation}
with $E_a\in\rest{\Z_p}^\times$, so $N_a\in\rest{\Z_p}$. We call
$\omega_a=2n_a-c_a-1$ the \emph{base cost} of the class $a$.

\subsection{Two lower bounds}\label{subsec:two-bounds}

\begin{proposition}[jet bound]\label{prop:jet}
$\vG p(\Delta)\ge\mathcal J_p$, where $\mathcal J_p$ is the minimum of
$\sum_a\sum_{j<t_a}(\omega_a+2j)$ over all $(t_a)_{a}$ with $t_a\ge0$ and $\sum_at_a=h$.
\end{proposition}

\begin{proof}
For polynomials $f$, $g$ of degree less than $h$ write $f(a+pz)=\sum_{i<h}f^{(a)}_ip^iz^i$; for
$f=u^r$ the Taylor coefficients $f^{(a)}_i=\binom ri a^{r-i}$ are integers. By
\cref{lem:distribution} and~\eqref{eq:local-weight},
\[
  \mu_X(\varphi u^ru^s)=\sum_{a=1}^p\sum_{i,j<h}(u^r)^{(a)}_i(u^s)^{(a)}_j\,\Lambda^{(a)}_{ij},
  \qquad
  \Lambda^{(a)}_{ij}=w^ap^{\,i+j-1+2n_a-c_a}\,\ell_a\Bigl(\frac{N_a(z)z^{i+j}}{P_a(z)}\Bigr),
\]
and $\vG p(\Lambda^{(a)}_{ij})\ge i+j+\omega_a$ by \cref{lem:local-integrality}(b). So the matrix of
$\Delta$ is $\mathbf J^{\mathsf T}\Lambda\mathbf J$ with $\mathbf J$ integral (rows indexed by
pairs $(a,i)$) and $\Lambda=\diag(\Lambda^{(a)})$ block diagonal. In~\eqref{eq:cauchy-binet} a
nonzero term selects $t_a$ rows and $t_a$ columns in each block, with distinct orders $i$ and
distinct orders $j$; each term of the Leibniz expansion of such a minor has valuation at least
$t_a\omega_a+2\binom{t_a}2=\sum_{j<t_a}(\omega_a+2j)$.
\end{proof}

Let $\epsilon_a=1$ if $\Pi_a$ contains an index $k$ with $2k+1<p$, and $\epsilon_a=0$
otherwise. Two elements of $\Pi_a$ differ by a multiple of $p$, so at most one index of $\Pi_a$
is below $(p-1)/2$, and it is the smallest one.

\begin{proposition}[pole bound]\label{prop:raw}
$\vG p(\Delta)\ge\mathcal R_p$, where $\mathcal R_p$ is the minimum of
$\sum_a\bigl[\sum_{j<t_a}(\omega_a+2j)+2\epsilon_a\one_{t_a=c_a}\bigr]$ over all $(t_a)$ with
$0\le t_a\le c_a$ and $\sum_at_a\le h$.
\end{proposition}

\begin{proof}
Write $\varphi u^s=Q_s+\sum_kx_k^s\rho_k/(u-x_k)$, $\rho_k=x_kN_A(x_k)/D_K'(x_k)$. The quotient
$Q_s$ of $u^{s+1}N_A$ by the monic $D_K\in\Z_p[u]$ is integral, so the matrix $\mathbf P=
[\mu_X(Q_{i+j})]$ is integral, and the matrix of $\Delta$ is
$\mathbf P+\sum_k\gamma_k\vvec(x_k)\vvec(x_k)^{\mathsf T}$ with
$\gamma_k=\rho_k\mu_X(1/(u-x_k))$. For $k\in\Pi_a$: $v_p(N_A(x_k))\ge2n_a$; exactly $c_a-1$
factors $l-k$ of $D_K'(x_k)$ are divisible by $p$, each exactly once; and $\mu_X(1/(u-x_k))$ has
Gauss valuation at least $-2$, and at least $0$ if $2k+1<p$. Hence
$\vG p(\gamma_k)\ge\omega_a+2\cdot\one_{2k+1<p}$.

Write the matrix as $[I_h\mid V]\diag(\mathbf P,\Gamma)[I_h\mid V]^{\mathsf T}$ with
$V=[\vvec(x_k)]_k$ and $\Gamma=\diag(\gamma_k)$. A nonzero term of~\eqref{eq:cauchy-binet}
selects the same set $\mathcal K$ of poles on both sides, $\abs{\mathcal K}\le h$, and some columns of
$I_h$. Its middle factor has valuation at least $\sum_{k\in\mathcal K}\vG p(\gamma_k)$. Each outer
factor is, up to sign, a minor of $V$ with columns $\mathcal K$; as a polynomial in the nodes it is
alternating with integer coefficients, hence divisible by $\prod_{k<l\in\mathcal K}(x_l-x_k)$, and
the $t_a=\abs{\mathcal K\cap\Pi_a}$ nodes of each class contribute valuation $\binom{t_a}{2}$. If
$t_a=c_a$, the index with $2k+1<p$, when it exists, is selected.
\end{proof}

The normalising factor contributes, by Legendre's formula and $p^2>3h>2K>h$,
\begin{equation}\label{eq:vp-S}
  v_p(\nrm)=h\Bigl(\floor{\tfrac Kp}-2\floor{\tfrac Ap}\Bigr)-2\sum_{i<h}\floor{\tfrac ip}
  =h\Bigl(\floor{\tfrac Kp}-2\floor{\tfrac Ap}-2Q_p+\tfrac phQ_p(Q_p+1)\Bigr),
\end{equation}
with $Q_p=\floor{h/p}$. Hence, for ordinary $p\ge\delta h$ and $h$ large,
\begin{equation}\label{eq:vp-T-combined}
  \vG p(T)\ \ge\ \max\bigl(-2h,\ v_p(\nrm)+\max(\mathcal J_p,\mathcal R_p)\bigr).
\end{equation}

\subsection{The profile}\label{subsec:profile}

\emph{Heuristic.} For $p=xh$ the residue classes come in a bounded number of types, according to
the numbers of poles and zeros they contain and to whether their first pole is improved, and the
number of classes of each type is $h$ times the length of an interval. The minima
$\mathcal J_p$ and $\mathcal R_p$ are then $h$ times the values of two filling problems whose
capacities are these lengths. We use the filling problems in their dual form, which makes both the
comparison with $\mathcal J_p$, $\mathcal R_p$ and the convexity in $x$ immediate.

\begin{lemma}[filling problems]\label{lem:filling}
Let $(e_i)_{i\in I}$ be real costs with capacities $m_i\ge0$, finitely many below any given
bound, and let $M>0$. For every selection $0\le y_i\le m_i$ with $\sum_iy_i=M$ and every
$\lambda\in\R$,
\[
  \sum_iy_ie_i\ \ge\ \lambda M-\sum_im_i\pos{\lambda-e_i},
\]
and the same holds for $\sum_iy_i\le M$ if $\lambda\le0$. Assume $\sum_im_i\ge M$. Equality holds for the
greedy selection of mass $M$ (cheapest costs first) with $\lambda$ its largest cost; in the second
case, for the greedy selection restricted to the negative costs (of mass $\min(M,\sum_{e_i<0}m_i)$)
with $\lambda$ the minimum of its largest cost and $0$.
\end{lemma}

\begin{proof}
$\sum_iy_ie_i=\lambda\sum_iy_i-\sum_iy_i(\lambda-e_i)$ and
$y_i(\lambda-e_i)\le m_i\pos{\lambda-e_i}$; if $\sum y_i\le M$ and $\lambda\le0$, then
$\lambda\sum_iy_i\ge\lambda M$. For the greedy selection of mass $M$ every $i$ with $e_i<\lambda$ is used
fully and every $i$ with $e_i>\lambda$ not at all, so both inequalities are equalities. In the
second case, if the negative costs have total capacity at least $M$, the greedy selection of mass
$M$ uses only negative costs and $\lambda\le0$ is its largest cost, so the first case applies;
otherwise take $y_i=m_i$ for $e_i<0$ and $y_i=0$ otherwise, and $\lambda=0$: both sides equal
$\sum_{e_i<0}m_ie_i$.
\end{proof}

\emph{The model.} Let $\kappa=K/h$ and $\alpha=A/h$ be fixed, and let $x>0$. Put
$c_0=\floor{\kappa/x}$, $r_\kappa=\kappa-c_0x$, $n_0=\floor{\alpha/x}$, $r_\alpha=\alpha-n_0x$, and
for $t\in[0,x)$
\begin{gather*}
  c(t)=c_0+\one_{t<r_\kappa},\qquad n(t)=n_0+\one_{((x/2-t)\bmod x)<r_\alpha},\\
  \epsilon(t)=\one_{t<x/2}\,\one_{c(t)>0},\qquad \omega(t)=2n(t)-c(t)-1.
\end{gather*}
Here $t$ stands for $k_0/h$, where $k_0$ is the least pole index in a class (see the proof of
\cref{prop:profile}). For $\lambda\in\R$ put
\[
  F_\lambda(\omega)=\sum_{j\ge0}\pos{\lambda-\omega-2j},\qquad
  \widetilde F_\lambda(\omega,c,\epsilon)=\sum_{0\le j<c}\pos{\lambda-\omega-2j-2\epsilon\one_{j=c-1}},
\]
and define
\begin{align*}
  J(x)&=\sup_{\lambda\in\R}\Bigl[\lambda-\int_0^xF_\lambda(\omega(t))\diff t\Bigr],\\
  R(x)&=\sup_{\lambda\le0}\Bigl[\lambda-\int_0^x\widetilde F_\lambda\bigl(\omega(t),c(t),
  \epsilon(t)\bigr)\diff t\Bigr],\\
  N(x)&=\floor{\kappa/x}-2\floor{\alpha/x}-2\floor{1/x}+x\floor{1/x}\bigl(\floor{1/x}+1\bigr),\\
  B(x)&=\max\bigl(-2,\ N(x)+\max(J(x),R(x))\bigr).
\end{align*}
The slots of the model take finitely many cost values $\omega(t)+2j$ below any bound, since
$\omega$ takes finitely many values; grouping the slots by cost value turns the continuous problems
into problems of the form of \cref{lem:filling}, with capacities the total length of $t$ with the
given cost. By \cref{lem:filling}, $J(x)$ is the least cost of filling mass $1$ with capacity $\diff t$ at each
cost $\omega(t)+2j$ (the \emph{jet slots}), and $R(x)$ the least cost of filling mass at most $1$
with capacity $\diff t$ at each cost $\omega(t)+2j+2\epsilon(t)\one_{j=c(t)-1}$, $j<c(t)$ (the
\emph{pole slots}); this primal form is the one evaluated in Appendix~\ref{app:profile-certificate}.

\begin{proposition}\label{prop:profile}
For every $\delta>0$ there is $C_\delta$ such that for all large $h$ in the construction and all
ordinary primes $\delta h\le p\le2K$, $\vG p(T)\ge h\,B(p/h)-C_\delta$.
\end{proposition}

\begin{proof}
Let $x=p/h\in[\delta,2\kappa]$. By~\eqref{eq:vp-S}, $v_p(\nrm)=hN(x)$ exactly, so
by~\eqref{eq:vp-T-combined} it suffices to show that the minima $\mathcal J_p$ and $\mathcal R_p$
of \cref{prop:jet,prop:raw} satisfy $\mathcal J_p\ge hJ(x)-C_\delta$ and
$\mathcal R_p\ge hR(x)-C_\delta$.

\emph{Exact class data.} For $0\le k_0<p$ let $x_{k_0}=-k_0-\frac12$ (a pole only if $k_0<K$).
These $p$ points lie in distinct classes modulo $p$, so the classes are indexed by $k_0$; the class
of $k_0$ contains the poles $x_k$ with $k\equiv k_0\pmod p$, $k<K$. Since $K/p=\kappa/x$ and
$K-p\floor{K/p}=hr_\kappa$, their number is $c(k_0/h)$ (for $K<p\le2K$ one has $c_0=0$,
$r_\kappa=\kappa$ and the classes with $k_0\ge K$ contain no pole, in accordance with
$c(t)=0$ for $t\ge r_\kappa$). Its first pole has
$2k_0+1<p$ if and only if $k_0\le(p-3)/2$, which agrees with $\one_{k_0/h<x/2}$ except for
$k_0=(p-1)/2$. Its zeros are the $j\in[1,A]$ with $j\equiv x_{k_0}\equiv\frac{p-1}2-k_0$; if
$\varrho\in\{1,\dots,p\}$ denotes this residue, their number is $n_0+\one_{\varrho\le hr_\alpha}$,
while $h\bigl((x/2-k_0/h)\bmod x\bigr)$ equals $\varrho+\frac12$ for $\varrho<p$ and $\frac12$ for
$\varrho=p$. Since $\varrho$ and $hr_\alpha$ are integers, the two indicators agree unless
$\varrho=hr_\alpha$ or $\varrho=p$. Altogether, at most three classes have data
$(c_a,n_a,\epsilon_a)$ different from $(c,n,\epsilon)(k_0/h)$.

\emph{Counting.} The model data are constant on each of the at most four intervals (open, closed or
half-open) into which the three interior points $x/2$, $r_\kappa$ and $(x/2-r_\alpha)\bmod x$ cut
$[0,x)$ (the set where $n(t)=n_0+1$ is $(s,x/2]$ cyclically, closed on the right), and an interval
of length $v-u$ contains between $h(v-u)-1$ and $h(v-u)+1$ of the points $k_0/h$; this gives an error of at most
$1$ per piece. Hence, for every $\lambda$,
\[
  \sum_aF_\lambda(\omega_a)\le h\int_0^xF_\lambda(\omega(t))\diff t+8\max_\omega F_\lambda(\omega),
\]
the maximum over $\omega\in[-c_0-2,\,2n_0+1]$, the range of all base costs; similarly for
$\widetilde F_\lambda$.

\emph{Comparison.} $\mathcal J_p$ is the least sum of $h$ elements of the multiset
$\{\omega_a+2j\}$, so \cref{lem:filling} gives $\mathcal J_p\ge h\lambda-\sum_aF_\lambda(\omega_a)$
for every $\lambda$. The supremum defining $J(x)$ is attained at the largest cost $\lambda^*$ used
by the greedy filling; as the levels $j<\lceil1/x\rceil$ have total capacity at least $1$,
$\lambda^*\le2n_0+1+2\lceil1/x\rceil$. With $\lambda=\lambda^*$ we obtain $\mathcal J_p\ge
hJ(x)-8\max_\omega F_{\lambda^*}(\omega)$, and $F_{\lambda^*}(\omega)\le(\lambda^*+c_0+4)^2/4$ is
bounded in terms of $\delta$, because $c_0\le\kappa/\delta$, $n_0\le\alpha/\delta$ and
$1/x\le1/\delta$. In the same way $\mathcal R_p$ is the least sum of at most $h$ elements of the
multiset of pole slots, and \cref{lem:filling} applies with $\lambda\le0$. The supremum over
$\lambda\le0$ defining $R(x)$ is attained, since the bracket is concave and piecewise affine in
$\lambda$ with slope $1$ for $\lambda<\min\omega-2$ (all terms vanish there), so a maximiser
$\lambda^*\le0$ exists. For $\lambda\le0$ and $\omega\ge-c_0-2$ each of the at most $c_0+1$
terms of $\widetilde F_\lambda$ is at most $c_0+2$, hence
$\widetilde F_\lambda\le(c_0+1)(c_0+2)$, and $\mathcal R_p\ge hR(x)-8(c_0+1)(c_0+2)$.
\end{proof}

\begin{lemma}[piecewise convexity]\label{lem:convexity}
Let $\mathfrak B$ be the set of points $1/j$, $\kappa/j$, $2\kappa/j$, $\alpha/j$, $2\alpha/j$,
$2(\kappa+\alpha)/j$, $j\ge1$ (only the first three if $\alpha=0$). On every open interval of
$(0,\infty)\setminus\mathfrak B$ the function $B$ is convex. Consequently
$\int_a^bB\ge(b-a)B(\frac{a+b}2)$ for consecutive points $a<b$ of $\mathfrak B\cup\{\delta\}$.
\end{lemma}

\begin{proof}
Let $U$ be an open interval disjoint from $\mathfrak B$. The integers $\floor{1/x}$, $c_0$, $n_0$
change only at the points $1/j$, $\kappa/j$, $\alpha/j$, so they are constant on $U$; hence $N$ is
affine on $U$, and $x/2$, $x$ and $r_\kappa=\kappa-c_0x$ are affine. The point
$s=(x/2-r_\alpha)\bmod x$ equals $x/2-r_\alpha$ or $x/2-r_\alpha+x$ according as
$r_\alpha\le x/2$ or not, and $r_\alpha=x/2$ means $x=2\alpha/(2n_0+1)\in\mathfrak B$, so $s$ is
affine on $U$ too. Two of the cut points $0$, $x/2$, $r_\kappa$, $s$, $x$ coincide only if
$r_\kappa=0$ ($x=\kappa/c_0$), $r_\kappa=x/2$ ($x=2\kappa/(2c_0+1)$), $s=0$ ($r_\alpha=x/2$),
$s=x/2$ ($r_\alpha=0$, $x=\alpha/n_0$) or $s=r_\kappa$ (then $\kappa+\alpha\in x/2+x\Z$, so
$x=2(\kappa+\alpha)/(2m+1)$); the remaining coincidences ($r_\kappa=x$, $s=x$, $x/2=0$) are
impossible. All these values lie in $\mathfrak B$, so on $U$ the cut points keep their order, the
model data $(c,n,\epsilon)$ are constant on each piece, and the lengths of the pieces are affine
in $x$. For fixed $\lambda$ the functions $x\mapsto\lambda-\int_0^xF_\lambda(\omega(t))\diff t$
and $x\mapsto\lambda-\int_0^x\widetilde F_\lambda\diff t$ are therefore affine on $U$, so $J$ and $R$, as
suprema of affine functions, are convex. They are finite: $R\le0$ and
$J\le2n_0+1+2\lceil1/x\rceil$. Hence $B$ is convex and bounded on $U$. A convex function on $(a,b)$ lies above
its supporting line at $\frac{a+b}2$, which gives the integral inequality.
\end{proof}

\subsection{The ordinary contribution}\label{subsec:ordinary-contribution}

\begin{theorem}\label{thm:ordinary}
For every $\delta>0$,
\[
  \liminf_{h\to\infty}\frac1{h^2}\sum_{p\ \mathrm{ordinary}}\vG p(T)\log p\ \ge\
  \int_\delta^{2\kappa}B(x)\diff x-2\delta .
\]
With $\delta=1/2000$ the right-hand side is at least $\Iprof_\Gen=-2.113164137$ for
$(\kappa,\alpha)=(7/6,0)$ and at least $\Iprof_\Band=-1.987089582$ for $(49/40,1/20)$.
\end{theorem}

\begin{proof}
Split the ordinary primes into $p<\delta h$, $\delta h\le p\le2K$ and $p>2K$. The last range
contributes a nonnegative amount by \cref{lem:large-primes}. For the first, by
\cref{prop:uniform},
\[
  \sum_{p<\delta h}2hL_p\log p\le2h\,\theta(\delta h)+2h\sqrt{3h}\log(3h)=2\delta h^2+o(h^2),
\]
because $L_p=1$ for $\sqrt{3h}<p<\delta h$, with $\theta(x)=\sum_{p\le x}\log p\sim x$. In the middle
range, \cref{prop:profile} and $\sum_{p\le2K}\log p=O(h)$ reduce the claim to
$\sum_{\delta h\le p\le2K}B(p/h)\log p\ge h\int_\delta^{2\kappa}B-o(h)$. The function $B$ is bounded
on $[\delta,2\kappa]$ and continuous outside the finite set $\mathfrak B\cap[\delta,2\kappa]$,
hence Riemann integrable; approximating it from below by step functions and using
$\theta(\beta h)-\theta(\beta'h)=(\beta-\beta')h+o(h)$ for fixed $\beta'<\beta$ (the prime number
theorem) gives the claim. The finitely many excluded primes do not affect the limit. The numerical
values are lower bounds for the midpoint sums of \cref{lem:convexity} over the
$5999$ (respectively $11798$) intervals cut out by $\mathfrak B$ in $[\delta,2\kappa]$, computed
exactly in rational arithmetic (Appendix~\ref{app:profile-certificate}).
\end{proof}

\begin{remark}\label{rem:profile-k1}
For $(\kappa,\alpha)=(1,0)$ the same computation gives $\int_\delta^2B-2\delta=-1.7169166\ldots
=\frac12-\frac{133}{60}-\frac\delta2$. Since $\log\nrm=\frac12h^2+O(h\log h)$ when $\kappa=1$,
this is Krisman's generic constant $\frac{133}{60}$ \cite[Thm.~4.5]{Krisman2026chi2} up to the
cut-off: at $\kappa=1$ the profile reproduces his class bound, and the gain for $\Gen$ does not
come from a sharper estimate at ordinary primes. What the profile adds is that it is defined for
every $\kappa$ and $\alpha$.
\end{remark}
